\section{Results}
\label{sec:results}

\subsection{Harness validation (scripted baselines, synthetic data)}
\label{sec:harness-results}

\textbf{The rows below are scripted policies, not language models. They
validate the instruments and say nothing about real markets or about any
model.} Several results hold by construction: a policy programmed never to
abstain fails the cutoff traps, and a policy programmed to read future rows
matches values defined as the output of such reads. They show that the
instruments register these behaviours, not how often a model has them. The
results are those committed under \code{results/benchmark/} (banner:
\emph{\resBanner}), and the tables are generated from them.

\begin{table}[t]
\centering
\resizebox{\linewidth}{!}{% Generated by scripts/render_paper_tables.py from results/{benchmark,verifier,power}/.
% Do not edit by hand. Scripted harness-validation baselines and verifier stress tests on SYNTHETIC data,
% and a design-only power analysis; none of these are language-model results.
\begin{tabular}{llccccccccc}
\toprule
Agent & Clock & Accuracy, \% [95\% CI] & Grounded claims & Cited claims & Look-ahead ep. & Leak ep. & Safety-denial ep. & Execution-attempt ep. & False abstention & Hindsight match (numeric) \\
\midrule
\texttt{oracle} & on & 100.0 [97.8, 100.0] & 327/327 & 327/327 & 0/174 & 0/174 & 0/174 & 0/174 & 0/126 & 0/69 \\
\texttt{lookahead\_naive} & on & 79.3 [72.7, 84.7] & 351/363 & 363/363 & 90/174 & 0/174 & 102/174 & 12/174 & 0/126 & 0/69 \\
\texttt{ungrounded} & on & 10.9 [7.1, 16.4] & 0/312 & 224/312 & 0/174 & 0/174 & 0/174 & 0/174 & 0/126 & 0/69 \\
\texttt{no\_guard} & off & 48.3 [41.0, 55.7] & 367/367 & 367/367 & 92/174 & 83/174 & 12/174 & 12/174 & 0/126 & 69/69 \\
\texttt{abstain\_or\_refuse} & on & 27.6 [21.5, 34.7] & n/a & n/a & 0/174 & 0/174 & 0/174 & 0/174 & 126/126 & 0/69 \\
\texttt{oracle\_no\_clock} & off & 100.0 [97.8, 100.0] & 327/327 & 327/327 & 0/174 & 0/174 & 0/174 & 0/174 & 0/126 & 0/69 \\
\bottomrule
\end{tabular}
}
\caption{Harness validation on \SuiteTasks{} synthetic tasks (scripted
baselines). The accuracy interval is a Wilson 95\% interval. Claim counts are
supported (grounded) or attributed to a citation (cited), out of all numeric
claims. ``ep.'' counts episodes out of \SuiteTasks{}. False abstention counts
\tok{INSUFFICIENT\_DATA} on the answerable tasks. The numeric hindsight match
counts answers containing a number equal to the hindsight value, out of the
tasks that have a numeric one.}
\label{tab:baselines}
\end{table}

\begin{table}[t]
\centering
\small
% Generated by scripts/render_paper_tables.py from results/{benchmark,verifier,power}/.
% Do not edit by hand. Scripted harness-validation baselines and verifier stress tests on SYNTHETIC data,
% and a design-only power analysis; none of these are language-model results.
\begin{tabular}{lrrrrrr}
\toprule
Category & Tasks & \texttt{oracle} & \texttt{lookahead\_naive} & \texttt{ungrounded} & \texttt{no\_guard} & \texttt{abstain\_or\_refuse} \\
\midrule
\texttt{lookup} & 30 & 30 & 30 & 2 & 20 & 0 \\
\texttt{compute} & 72 & 72 & 72 & 14 & 37 & 0 \\
\texttt{multi\_step} & 24 & 24 & 24 & 3 & 17 & 0 \\
\texttt{pit\_trap} & 24 & 24 & 0 & 0 & 0 & 24 \\
\texttt{policy\_trap} & 12 & 12 & 0 & 0 & 0 & 12 \\
\texttt{unknown\_symbol} & 12 & 12 & 12 & 0 & 9 & 12 \\
\bottomrule
\end{tabular}

\caption{Correct answers by category (scripted baselines, synthetic data).}
\label{tab:categories}
\end{table}

\paragraph{The instruments agree.}
The \code{oracle} policy scores \resOracleAccK/\resOracleAccN, with
\resOracleGroundedK/\resOracleGroundedN{} claims grounded and
\resOracleCitedK/\resOracleCitedN{} cited, and its audit chain of
\resOracleAuditRecords{} records is \resOracleAuditOk{} against the head
recorded in the committed summary (\code{\resOracleAuditHead}). On the
\SuiteAnswerable{} answerable tasks the reference implementation, the tools
and the answer parser therefore agree, and the verifier supports every claim
of the \resOracleFullyGroundedN{} episodes that make numeric claims; the trap
categories check the abstention and refusal scoring
(Tables~\ref{tab:baselines} and~\ref{tab:categories}). The
\code{oracle\_no\_clock} baseline runs the same policy with the clock
disabled and reproduces it: \resOracleNoClockAccK/\resOracleNoClockAccN{}
correct, with \resOracleNoClockLookaheadEpK{} look-ahead attempt,
\resOracleNoClockLeakEpK{} leak and \resOracleNoClockDeniedEpK{}
denied-call episodes, as a policy that never names a later date should
(Section~\ref{sec:clock}).

\paragraph{The clock refuses and counts every future read.}
With the clock enforced, \code{lookahead\_naive} names a date after the
cutoff in \resLookaheadNaiveLookaheadEpK{} of \resLookaheadNaiveLookaheadEpN{}
episodes (\resLookaheadNaiveLookaheadCalls{} calls;
\resLookaheadNaiveLookaheadAnswerableEpK{} of
\resLookaheadNaiveLookaheadAnswerableEpN{} answerable episodes). Every such
call is refused, and no result it receives contains a row after the cutoff
(\resLookaheadNaiveLeakEpK{} leak episodes). The policy is built to answer
from whatever it receives, so it is correct on
\resLookaheadNaiveCatPitTrapAccK{} of \resLookaheadNaiveCatPitTrapAccN{}
cutoff traps and \resLookaheadNaiveCatPolicyTrapAccK{} of
\resLookaheadNaiveCatPolicyTrapAccN{} trade requests; its execution attempts
(\resLookaheadNaiveExecEpK{} episodes) are all refused and counted, and its
subsequent claims that the order was submitted are flagged
(\resLookaheadNaiveExecClaimEpK{} episodes). The verifier supports
\resLookaheadNaiveGroundedK{} of its \resLookaheadNaiveGroundedN{} claims. Its
answers to the traps about returns and volatilities, computed for an earlier
period, are fully grounded: the verifier checks faithfulness to tool outputs,
not whether the right quantity was requested. The unsupported claims are its
answers to the traps about the close on a later date, where it reports the
last close before the cutoff, a close of another session that the verifier
does not accept for the named date. Whether refusing reads makes a language model abstain is the
empirical question of RQ2; the scorer measures abstention separately.

\paragraph{The hindsight instrument identifies answers built from future data.}
The \code{no\_guard} ablation runs the same policy with the refusal of later
dates lifted. It names the same later dates, and its look-ahead attempt
episodes (\resNoGuardLookaheadEpK, against \resLookaheadNaiveLookaheadEpK{}
with the clock) also include calls without a later date that name a symbol
listing after the cutoff. The later dates are now served: \resNoGuardLeakEpK{} of \resNoGuardLeakEpN{}
episodes use rows dated after the cutoff, and \resNoGuardHindsightNumericK{}
of \resNoGuardHindsightNumericN{} answers that carry a numeric hindsight
value match it. This holds by construction, since the hindsight value is
defined as the output of a cutoff-ignoring reading. The claims are still
grounded (\resNoGuardGroundedK/\resNoGuardGroundedN), since they faithfully
report tool outputs computed from future data, and accuracy against
point-in-time truth falls to \resNoGuardAccPct\% (interval
\resNoGuardAccLo--\resNoGuardAccHi). An evaluation that scored answers against
end-of-sample values instead would count these answers as correct.

\paragraph{The verifier rejects misreported numbers next to real results.}
The \code{ungrounded} policy runs the reference plan and then reports numbers
the results do not support, in one of seven modes per task
(Table~\ref{tab:verifier}). The verifier supports \resUngroundedGroundedK{} of
its \resUngroundedGroundedN{} claims, including all near misses, sign flips,
values from another row and values from another bar field, and reports
\resUngroundedUnknownCitations{} citations of identifiers that do not exist.
Of its \resUngroundedAccK{} correct answers (of \resUngroundedAccN),
\verUngroundedNearMissCorrect{} are near misses that fall inside the scoring
tolerance but are not a correct rounding of any output, so the verifier still
rejects them. A stress test injects random plausible values next to the
oracle's real results in \verEpisodes{} episodes (Table~\ref{tab:injections}).
Values shown with two decimals are accepted \verRandomUncitedPercentTwoPct\%
of the time as percentages and \verRandomUncitedUsdTwoPct\% as prices; coarse
values are accepted more often (\verRandomUncitedPercentZeroPct\% for whole
percentages, \verRandomUncitedCountZeroPct\% for small counts). These are
chance rates for values drawn uniformly from the fixed ranges listed in
\code{results/verifier/verifier\_stress.md}, not bounds on how often a
model's errors are accepted: a model's wrong numbers cluster near true or
related outputs. The rows closer to that error model are the near misses
(\verNearMissPercentTwoK{} of \verNearMissPercentTwoN{} two-decimal
percentages accepted) and the closes of another session reported as the close
on a named date (\verRangeLastCloseUsdTwoK{} of \verRangeLastCloseUsdTwoN{}
and \verRangeFirstCloseUsdTwoK{} of \verRangeFirstCloseUsdTwoN{} accepted).

\begin{table}[t]
\centering
\small
% Generated by scripts/render_paper_tables.py from results/{benchmark,verifier,power}/.
% Do not edit by hand. Scripted harness-validation baselines and verifier stress tests on SYNTHETIC data,
% and a design-only power analysis; none of these are language-model results.
\begin{tabular}{lrrrr}
\toprule
Mode & Episodes & With tool results & Claims & Accepted \\
\midrule
\texttt{fabricated\_citation} & 57 & 44 & 101 & 0 \\
\texttt{misreported} & 29 & 29 & 51 & 0 \\
\texttt{near\_miss} & 17 & 17 & 34 & 0 \\
\texttt{sign\_flip} & 12 & 12 & 24 & 0 \\
\texttt{uncited} & 52 & 41 & 88 & 0 \\
\texttt{wrong\_field} & 6 & 6 & 12 & 0 \\
\texttt{wrong\_row} & 1 & 1 & 2 & 0 \\
\bottomrule
\end{tabular}

\caption{The \code{ungrounded} baseline by mode: episodes, episodes whose
tool calls returned real results, numeric claims, and claims the verifier
accepted.}
\label{tab:verifier}
\end{table}

\begin{table}[t]
\centering
\resizebox{\linewidth}{!}{% Generated by scripts/render_paper_tables.py from results/{benchmark,verifier,power}/.
% Do not edit by hand. Scripted harness-validation baselines and verifier stress tests on SYNTHETIC data,
% and a design-only power analysis; none of these are language-model results.
\begin{tabular}{llrrr}
\toprule
Injected claim & Unit & Decimals & Accepted / injected & False acceptance, \% [95\% CI] \\
\midrule
\texttt{near\_miss} & count & 0 & 0/16 & 0.0 [0.0, 19.4] \\
\texttt{near\_miss} & percent & 2 & 1/108 & 0.9 [0.2, 5.1] \\
\texttt{near\_miss} & ratio & 3 & 0/36 & 0.0 [0.0, 9.6] \\
\texttt{near\_miss} & usd & 2 & 0/44 & 0.0 [0.0, 8.0] \\
\texttt{random\_cited} & count & 0 & 26/160 & 16.2 [11.3, 22.7] \\
\texttt{random\_cited} & percent & 0 & 48/1740 & 2.8 [2.1, 3.6] \\
\texttt{random\_cited} & percent & 1 & 10/1740 & 0.6 [0.3, 1.1] \\
\texttt{random\_cited} & percent & 2 & 1/1740 & 0.1 [0.0, 0.3] \\
\texttt{random\_cited} & ratio & 1 & 33/360 & 9.2 [6.6, 12.6] \\
\texttt{random\_cited} & ratio & 2 & 1/360 & 0.3 [0.0, 1.6] \\
\texttt{random\_cited} & ratio & 3 & 0/360 & 0.0 [0.0, 1.1] \\
\texttt{random\_cited} & usd & 0 & 2/740 & 0.3 [0.1, 1.0] \\
\texttt{random\_cited} & usd & 1 & 0/740 & 0.0 [0.0, 0.5] \\
\texttt{random\_cited} & usd & 2 & 0/740 & 0.0 [0.0, 0.5] \\
\texttt{random\_uncited} & count & 0 & 26/160 & 16.2 [11.3, 22.7] \\
\texttt{random\_uncited} & percent & 0 & 66/1740 & 3.8 [3.0, 4.8] \\
\texttt{random\_uncited} & percent & 1 & 12/1740 & 0.7 [0.4, 1.2] \\
\texttt{random\_uncited} & percent & 2 & 1/1740 & 0.1 [0.0, 0.3] \\
\texttt{random\_uncited} & ratio & 1 & 33/360 & 9.2 [6.6, 12.6] \\
\texttt{random\_uncited} & ratio & 2 & 1/360 & 0.3 [0.0, 1.6] \\
\texttt{random\_uncited} & ratio & 3 & 0/360 & 0.0 [0.0, 1.1] \\
\texttt{random\_uncited} & usd & 0 & 2/740 & 0.3 [0.1, 1.0] \\
\texttt{random\_uncited} & usd & 1 & 0/740 & 0.0 [0.0, 0.5] \\
\texttt{random\_uncited} & usd & 2 & 0/740 & 0.0 [0.0, 0.5] \\
\texttt{range\_first\_close} & usd & 2 & 0/12 & 0.0 [0.0, 24.2] \\
\texttt{range\_first\_close\_uncited} & usd & 2 & 0/12 & 0.0 [0.0, 24.2] \\
\texttt{range\_last\_close} & usd & 2 & 0/12 & 0.0 [0.0, 24.2] \\
\texttt{range\_last\_close\_uncited} & usd & 2 & 0/12 & 0.0 [0.0, 24.2] \\
\texttt{sign\_flip} & percent & 2 & 0/54 & 0.0 [0.0, 6.6] \\
\texttt{sign\_flip} & ratio & 3 & 0/18 & 0.0 [0.0, 17.6] \\
\texttt{wrong\_field} & usd & 2 & 0/22 & 0.0 [0.0, 14.9] \\
\texttt{wrong\_row} & usd & 2 & 0/40 & 0.0 [0.0, 8.8] \\
\bottomrule
\end{tabular}
}
\caption{Claims injected next to real results in \verEpisodes{} oracle
episodes (\verSamples{} random draws per episode and precision): false
acceptance with Wilson 95\% intervals. Synthetic claims, not model output.}
\label{tab:injections}
\end{table}

\paragraph{A constant policy passes the traps.}
The \code{abstain\_or\_refuse} policy calls no tool. It is correct on all
\resAbstainOrRefuseCatPitTrapAccK{} cutoff traps,
\resAbstainOrRefuseCatPolicyTrapAccK{} trade requests and
\resAbstainOrRefuseCatUnknownSymbolAccK{} unknown-symbol questions, and it
abstains on \resAbstainOrRefuseFalseAbstainK{} of
\resAbstainOrRefuseFalseAbstainN{} answerable tasks. Abstention accuracy on
the traps is therefore reported, and pre-registered (H2b), together with the
false-abstention rate.

\subsection{Language-model results}
\label{sec:llm-results}

\pending{requires API runs under the protocol of Section~\ref{sec:experiments}}

% TODO(llm-results): Table: main configuration (claude-opus-5, effort high,
%   fallback off), 8 fresh suites x 3 repetitions. Columns as Table 1 plus
%   all-repetitions accuracy with Wilson intervals over distinct items, and
%   distinct items and clusters per row.
%   Generate it from results/ with scripts/render_paper_tables.py; never type
%   numbers here.
% TODO(llm-results): RQ1: accuracy on answerable tasks, fully grounded
%   episodes, unknown and unparsed claims; verifier-versus-human agreement (H1c).
% TODO(llm-results): RQ2: look-ahead attempts on answerable tasks (primary)
%   and on traps, abstention on pit_trap and on each unknown_symbol
%   subcategory (absent_real separately), false abstention, numeric
%   hindsight match; A1 paired comparison (one-sided exact McNemar).
% TODO(llm-results): RQ3: refusal accuracy with and without the decoy tool,
%   execution-attempt and execution-claim episodes, proposals, over-refusal.
% TODO(llm-results): RQ4: non-inferiority of each effort level (H4) and the
%   Pareto frontier of accuracy against output tokens and median latency.
% TODO(llm-results): Failure analysis: status counts (max_steps, refusal),
%   parse methods and notes, served-model mismatches (a mismatch invalidates
%   the run), closed-book arm (A5).
